\documentclass[12pt, a4paper]{article}

\usepackage[margin=1.2in]{geometry}
\usepackage[utf8]{inputenc}
\usepackage[T1]{fontenc}
\usepackage{microtype}
\usepackage{amsmath, amssymb}
\usepackage{enumitem}
\usepackage{titlesec}
\usepackage{parskip}
\usepackage{hyperref}
\usepackage[round, authoryear]{natbib}

\hypersetup{
  colorlinks = true,
  linkcolor  = black,
  citecolor  = black,
  urlcolor   = black
}

\titleformat{\section}{\large\bfseries}{\thesection.}{0.5em}{}
\titleformat{\subsection}{\normalsize\bfseries}{\thesubsection}{0.5em}{}
\titleformat{\subsubsection}{\normalsize\itshape}{\thesubsubsection}{0.5em}{}

\setlength{\parindent}{0pt}
\setlength{\parskip}{0.75em}

\title{\textbf{Formal Mathematical Representation Discovery}\\[0.4em]
       \large A Research Thesis for AI-Native Mathematics}
\author{}
\date{}

% -----------------------------------------------------------------------
\begin{document}
\maketitle

\begin{abstract}
This thesis proposes \emph{formal mathematical representation discovery}: a
research program that treats Lean/mathlib as a machine-readable record of
mathematical abstraction, and builds AI systems that mine proof-state
transition graphs for reusable motifs, retrieve structurally analogous
theorems across domains, and detect missing abstractions through
graph-theoretic bottlenecks. The central claim is that the next major advance
in AI for mathematics will come not from better tactic prediction but from
learning how representations are discovered, compressed, and
transferred---a process that formal libraries now expose empirically for the
first time.
\end{abstract}

% -----------------------------------------------------------------------
\section{Central Thesis}

Modern formal mathematics libraries such as Lean/mathlib
\citep{mathlib2020} are not merely repositories of theorems and proofs.
They are large-scale, machine-readable records of how mathematical knowledge
is represented, compressed, transported, and reused. Their internal
structures---definitions, typeclasses, proof states, tactic traces,
dependency graphs, coercion paths, simplification rules, and theorem
statements---encode a computational shadow of mathematical abstraction
itself.

The central thesis of this research program is:

\begin{quote}
The next major advance in AI for mathematics will not come only from better
next-tactic prediction, larger language models, or brute-force proof search.
It will come from learning how mathematical representations are discovered,
reused, compressed, and transferred across domains.
\end{quote}

A formal mathematics system such as Lean gives us, for the first time, an
empirical substrate for studying this process. It exposes the topology of
proof search, the dependency geometry of mathematical concepts, the recurring
motifs of proof construction, and the places where existing representations
are insufficient. By mining these structures, we can build systems that
retrieve analogies, suggest intermediate lemmas, detect missing abstractions,
and eventually help invent new mathematical representations.

This research direction can be summarized as \textbf{formal mathematical
representation discovery}.

Its goal is not simply to prove more Lean theorems. Its goal is to
understand and automate the higher-level acts that make proof possible:
choosing the right object, introducing the right invariant, transporting
structure across domains, and compressing repeated proof patterns into
reusable abstractions.

% -----------------------------------------------------------------------
\section{Motivation}

Current AI theorem proving is largely framed as a local search problem. A
model sees a proof state and predicts a tactic, a lemma, or a proof term.
Recent systems such as AlphaProof \citep{alphaproof2025},
DeepSeek-Prover-V2 \citep{deepseekproverv2}, and HyperTree Proof Search
\citep{lample2022hypertree} have demonstrated that this framing, combined
with large-scale data and reinforcement learning, can achieve remarkable
results. Yet the framing is incomplete.

Human mathematical creativity rarely consists of local tactic selection
alone. Much of it consists of \emph{representation change}.

A difficult theorem often becomes tractable only after one introduces a new
object, defines a better interface, proves an intermediate lemma, recognizes
an analogy with another field, or reformulates the problem in a more natural
language. In mature mathematics, these representational moves are often more
important than the final proof script.

Formalized mathematics contains abundant evidence of this phenomenon. In
mathlib, many definitions and structures exist not merely because they are
mathematically natural in isolation, but because they make large families of
theorems reusable. A typeclass hierarchy, a bundled structure, a coercion
mechanism, a normal form, or a carefully shaped lemma can dramatically reduce
proof complexity across hundreds of downstream results.

This suggests a new view:

\begin{quote}
A mathematical abstraction is valuable when it compresses proof search and
increases theorem transportability.
\end{quote}

Lean makes this view measurable. We can ask:
\begin{itemize}[leftmargin=1.5em]
  \item Which structures reduce proof length and dependency depth?
  \item Which lemmas act as bottlenecks or bridges between domains?
  \item Which proof-state motifs recur across apparently unrelated areas?
  \item Which theorem statements have analogous relational structure despite
        different terminology?
  \item Where do proofs repeatedly suffer from the same representational
        friction?
  \item Can we detect missing abstractions before humans define them?
\end{itemize}

These questions move beyond theorem proving as local search and toward
theorem proving as representation engineering.

% -----------------------------------------------------------------------
\section{Related Work}

\subsection{AI Theorem Proving Systems}

The dominant paradigm in AI theorem proving frames proof search as a
sequential decision problem. AlphaProof \citep{alphaproof2025} combines a
language model with AlphaZero-style reinforcement learning inside Lean,
achieving silver-medal performance at the 2024 IMO. DeepSeek-Prover-V2
\citep{deepseekproverv2} uses RL-guided subgoal decomposition to achieve
88.9\% pass rate on MiniF2F. HyperTree Proof Search
\citep{lample2022hypertree} applies an online AlphaZero variant over and-or
proof trees. These systems advance the state of the art in automated proof
but operate entirely within the tactic-prediction paradigm this thesis moves
beyond.

\subsection{Premise Selection and Retrieval-Augmented Proving}

Premise selection---choosing which library lemmas to supply to a
prover---is the earliest form of library-aware proof automation. The hammer
tradition \citep{paulson2010sledgehammer, blanchette2016hammering} introduced
tools like Sledgehammer for Isabelle, with premise filters based on symbol
overlap \citep{meng2009mepo} and machine-learned classifiers
\citep{kuhlwein2013mash}. \citet{kaliszyk2014} demonstrated that ML-guided
premise selection at scale substantially increases automated proof coverage.
LeanDojo \citep{yang2023leandojo} extends this to Lean~4, providing
programmatic proof-state extraction and a retrieval-augmented prover
(ReProver). LeanSearch \citep{gao2024leansearch} provides semantic search
over Mathlib4 for informal natural-language queries. COPRA
\citep{thakur2023copra} uses GPT-4 with retrieved lemmas and stateful
backtracking.

This thesis's analogy retrieval component (\S\ref{sec:retrieval}) is related
to premise selection but differs in aim: rather than retrieving lemmas to
supply to a prover, it retrieves structurally analogous theorems as a source
of proof strategy, using multi-layer graph structure rather than text or
symbol overlap.

\subsection{Graph-Based Representations of Formal Libraries}
\label{sec:related-graphs}

Two bodies of work have applied graph analysis to formal proof libraries,
forming the closest technical predecessors to this thesis.

\citet{paliwal2020graph} and \citet{bansal2019holist} apply graph neural
networks to HOL Light formula graphs for tactic and premise prediction. These
systems use graph structure within individual theorem statements but do not
mine cross-theorem motifs or study abstraction at the library level.

\citet{li2026networkstructure} present a network analysis of Lean~4's
Mathlib itself, extracting its dependency structure into a multilayer graph
with over 308,000 declarations and millions of edges across thousands of
modules, and introduce graph decompositions distinguishing explicit from
compiler-synthesized edges. This is the closest existing work to the graph
extraction proposed in \S\ref{sec:extraction}. The key difference is scope
and use: \citeauthor{li2026networkstructure} study Mathlib's network structure
descriptively; this thesis proposes to use graph structure operationally for
motif retrieval, analogy search, and abstraction detection.

\citet{huch2022} applies the same network-analysis approach to the Isabelle
Archive of Formal Proofs (1.8M nodes, 2.8M edges). Like
\citeauthor{li2026networkstructure}, this work examines static dependency
centrality and does not pursue motif mining, analogy retrieval, or abstraction
detection.

\citet{blanchette2015mining} mine the AFP empirically for proof statistics:
proof size, tactic usage distributions, dependency structure, and proof style
evolution. This is the most direct precursor to the motif mining proposal in
\S\ref{sec:motif-mining}. The key distinction is that
\citeauthor{blanchette2015mining} report descriptive statistics over tactic
sequences; this thesis proposes to extract structured motifs from proof-state
transition graphs and use them operationally for retrieval and abstraction
suggestion.

PACT \citep{han2021pact} extracts self-supervised training tasks from Lean
proof terms. LeanDojo \citep{yang2023leandojo} extracts proof states,
premises, and tactic traces from Lean~4 at scale. The trace collection layer
(\S\ref{sec:arch-layer1}) builds on and extends this infrastructure.

\subsection{Proof-Pattern Mining and Tactic Discovery}

\paragraph{Proof-pattern mining.}
The most direct conceptual ancestors of this thesis are the ML4PG family of
systems for Coq/SSReflect. \citet{komendantskaya2012ml4pg} introduce ML4PG,
a machine learning plugin that clusters proof scripts by statistical features
extracted from proof states and tactic sequences, identifying similarity
across proof developments. \citet{heras2014recycling} demonstrate that ML4PG
can mine electronic Coq proof libraries and suggest proof strategies by
analogy. This is the clearest ancestor of Hypothesis~1 and
\S\ref{sec:motif-mining}. The key distinctions from this thesis are:
ML4PG uses shallow statistical features over tactic sequences rather than
structured proof-state transition graphs; it operates on Coq/SSReflect at a
scale far smaller than Lean/mathlib; and it does not address theorem analogy
retrieval or missing abstraction detection.

\citet{freitas2014} identify named, reusable ``proof patterns'' for formal
verification---recurring solutions to common proof obligations---through
manual expert identification; this thesis proposes automated discovery from
data at scale.

\paragraph{Tactic discovery.}
TacMiner \citep{xin2025tacminer} is the closest prior work to Project~1
(\S\ref{sec:agenda}). It builds Tactic Dependence Graphs over Lean~4 proofs
to identify reusable proof strategies across multiple proofs, discovers new
tactic library entries, and refactors proofs into more modular forms---reporting
3$\times$ as many learned tactics as Peano and a 26\% proof-size reduction.
The distinction from this thesis is one of scope: TacMiner targets tactic
library compression as an end in itself; this thesis uses motif structure as a
stepping stone toward theorem analogy retrieval, missing abstraction detection,
and representation discovery at library scale.

\subsection{Mathematical Analogy in Automated Reasoning}

Work on proof analogy in automated reasoning dates to the 1980s.
\citet{boydelatour1987} propose second-order pattern matching to transform
known proof terms into candidates for analogous propositions.
\citet{melis1994} use derivational analogy at the level of proof plans to
guide automated provers by replaying known analogous proofs. This thesis
pursues analogy at a different level of abstraction---structural graph
similarity across a large formal library---and aims at retrieval rather than
direct proof transfer.

The cognitive science foundation for relational analogy is structure-mapping
theory \citep{gentner1983, falkenhainer1989sme}, which holds that analogical
reasoning aligns relational structure between domains rather than surface
features. The multi-layer graph retrieval proposed in \S\ref{sec:retrieval}
can be viewed as a large-scale, data-driven realization of structure-mapping
principles in a formal mathematical setting.

\subsection{Theory Exploration and Conjecture Generation}

Theory exploration systems automatically discover lemmas in a given formal
theory development. Hipster \citep{johansson2014hipster} integrates theory
exploration into Isabelle/HOL using QuickSpec-style random testing to generate
equational conjectures, then attempts to prove or prune them. LeanConjecturer
\citep{onda2025leanconjecturer} generates university-level mathematical
conjectures for Lean~4 using LLMs, producing 12,289 conjectures from 40
Mathlib seed files.

Both are conceptually related to \S\ref{sec:agenda} (Missing Abstraction
Detection) in that both attempt to identify mathematical statements ``missing''
from a library. This thesis's approach differs by using proof-state transition
graphs and dependency topology to detect representation bottlenecks (repeated
motifs, long coercion chains, high-branching bottleneck states) and evaluating
suggestions by proof compression rather than by provability alone.

\subsection{Library Learning and Abstraction Discovery}

The program synthesis community has studied abstraction discovery as a
library learning problem. DreamCoder \citep{ellis2021dreamcoder} jointly
learns reusable library abstractions and a neural search policy through
wake-sleep iteration. Stitch \citep{bowers2023stitch} provides a
significantly more efficient top-down abstraction algorithm using program
compression ratio as the primary metric of abstraction quality. LAPS
\citep{wong2021laps} extends DreamCoder with natural language annotations to
guide library learning.

These systems are the program synthesis analogues of what this thesis
proposes for formal mathematics. The key differences are: (i)~the objects
are types and propositions rather than input-output examples; (ii)~correctness
is guaranteed by Lean's type checker rather than test suites; (iii)~the
abstraction space includes typeclasses, coercions, and universal properties,
which have no direct program synthesis analogue; and (iv)~the corpus is a
mature human-curated library rather than a set of synthetic tasks.

\subsection{Perspectives on Mathematical Abstraction}

\citet{mitchell2021} reviews AI approaches to abstraction and
analogy-making, concluding that no current system achieves human-level
analogical generalization and proposing evaluation criteria.
\citet{zhang2023cognitive} examine AI for mathematics from a cognitive
science perspective, arguing that current systems lack the
representation-change capabilities that make human mathematical reasoning
powerful---a diagnosis this thesis takes as its starting point.

% -----------------------------------------------------------------------
\section{Core Hypotheses}

This research program rests on five concrete hypotheses.

\subsection{Hypothesis 1: Proof Traces Contain Recurring Cross-Domain Motifs}

Lean proofs are not arbitrary tactic sequences. They contain recurring
structural patterns: introducing witnesses, decomposing conjunctions,
applying extensionality, transporting along equivalences, normalizing
coercions, invoking universal properties, reducing goals by induction, or
closing branches through simplification.

These motifs often recur across different mathematical domains. Prior
empirical work on the Isabelle AFP \citep{blanchette2015mining} and manual
pattern identification \citep{freitas2014} support the existence of such
patterns; this thesis proposes to discover them automatically from
proof-state transition graphs in Lean/mathlib.

\subsection{Hypothesis 2: Theorem Analogy Can Be Retrieved Through Relational
Structure}

Many important mathematical analogies are not lexical. Two theorems may
share little vocabulary but have similar dependency neighborhoods, type
structures, proof skeletons, or relational patterns.

Structure-mapping theory \citep{gentner1983, falkenhainer1989sme} predicts
that the relevant similarity is relational, not featural. A retrieval system
based on multi-layer formal graphs may find relational analogies that
text-based systems such as LeanSearch \citep{gao2024leansearch} miss.

\subsection{Hypothesis 3: Missing Abstractions Leave Measurable Traces}

When a library lacks the right abstraction, users compensate through repeated
proof patterns, duplicated lemmas, long coercion chains, brittle rewrite
sequences, and similar theorem statements spread across namespaces. These
phenomena should be detectable, turning abstraction discovery into a
mathematical refactoring problem analogous to code clone detection
\citep{roycordy2007, koschke2007}.

\subsection{Hypothesis 4: Abstraction Quality Can Be Measured by Compression}

A good abstraction reduces the cost of future proof. Quantitative measures
of abstraction value include: proof length reduction, dependency reuse,
decrease in branching factor, reduction of repeated proof motifs, increased
theorem transportability, and decrease in coercion and normalization burden.
This compression-as-quality view is operationalized in program synthesis by
\citet{bowers2023stitch}; this thesis applies it in the richer setting of
dependent type theory.

\subsection{Hypothesis 5: Proof-State Topology Predicts Difficulty and
Useful Intermediate Lemmas}

A proof is a trajectory through a space of proof states. Some intermediate
states act as bottlenecks: once reached, many downstream goals become easy.
These bottleneck states may correspond to useful intermediate lemmas,
invariants, or representation changes, and their geometry may guide lemma
invention and proof planning.

% -----------------------------------------------------------------------
\section{Research Agenda}
\label{sec:agenda}

The agenda has four mutually reinforcing components.

\subsection{Lean Graph Extraction and Normalization}
\label{sec:extraction}

The first task is to construct a graph-based representation of Lean/mathlib
that is rich enough to capture mathematical structure but normalized enough
to avoid irrelevant syntactic noise.

Existing tools provide a starting point. LeanDojo \citep{yang2023leandojo}
extracts proof states, premise annotations, and tactic traces from Lean~4
repositories at scale. PACT \citep{han2021pact} demonstrates that rich
training data can be extracted from Lean proof terms. This project extends
such infrastructure to produce a multi-layer graph intermediate
representation.

The system should extract multiple graph layers:

\subsubsection{Declaration dependency graph}
Nodes are definitions, theorems, instances, structures, and lemmas. Edges
represent dependency, invocation, unfolding, instance use, or theorem
application.

\subsubsection{Expression graph}
Theorem statements, definitions, and proof terms are represented as typed
expression graphs capturing binders, applications, constants, equalities,
propositions, functions, inductive types, and typeclass constraints.

\subsubsection{Typeclass and coercion graph}
This graph records how mathematical structures are connected through
instances, coercions, inheritance, bundled structures, and canonical
constructions.

\subsubsection{Proof-state transition graph}
For tactic proofs, each tactic transforms one or more goals into successor
goals, inducing a transition graph over proof states. The graph records local
context, target type, generated subgoals, tactic class, and closure
mechanism.

\subsubsection{Simplification and rewriting graph}
\texttt{simp}, \texttt{rw}, normalization, and definitional equality hide
substantial reasoning. The system exposes, where possible, the rewrite rules
and normalization paths used to close goals.

The central design problem is normalization. Raw Lean expressions contain
variable names, elaboration artifacts, syntactic choices, universe levels,
implicit arguments, coercions, and proof-irrelevant details. The system must
quotient out accidental differences while preserving mathematically meaningful
structure.

\subsection{Proof Motif Mining}
\label{sec:motif-mining}

Once proof-state transition graphs are available, the next project is to
discover recurring proof motifs.

A proof motif is a repeated local or global pattern in proof-state evolution.
It is not merely a tactic sequence---it is a structural transformation
pattern over goals, hypotheses, and generated subgoals.

The ML4PG family of systems \citep{komendantskaya2012ml4pg, heras2014recycling}
and empirical analysis of the Isabelle AFP \citep{blanchette2015mining} and
manually curated proof patterns \citep{freitas2014} all provide prior evidence
that recurring structures exist. The open question is whether they can be
discovered automatically from proof-state transition graphs in a form useful
for retrieval and proof planning, at the scale of Lean/mathlib.

The project would mine frequent subgraphs from proof-state transition graphs
and cluster them by structural similarity, producing a library of proof
motifs.

\paragraph{Evaluation.}
\begin{enumerate}[leftmargin=1.5em]
  \item \textit{Cross-domain recurrence}: Do motifs appear across multiple
        namespaces and mathematical areas?
  \item \textit{Predictive value}: Does motif identity help predict the next
        tactic or proof step?
  \item \textit{Retrieval value}: Can motif-aware retrieval find useful
        analogous proofs given a proof state?
  \item \textit{Compression value}: Can motifs summarize many proofs more
        compactly than tactic sequences?
  \item \textit{Human interpretability}: Do mathematicians or Lean users
        recognize mined motifs as meaningful strategies?
\end{enumerate}

\subsection{Theorem Analogy Retrieval}
\label{sec:retrieval}

The next component is a retrieval system that finds structurally analogous
theorems across mathlib. Given a theorem statement, proof state, or partial
proof, the system returns other theorems similar not merely in text but in
mathematical structure---motivated by structure-mapping theory
\citep{gentner1983, falkenhainer1989sme} and prior work on proof-plan
analogy \citep{melis1994}.

Similarity is computed using a combination of theorem statement expression
graphs, dependency neighborhoods, typeclass neighborhoods, proof motif
signatures, proof-state transition summaries, namespace distance, shared
abstraction usage, graph embeddings, and exact or approximate subgraph
matching.

\paragraph{Evaluation.}
\begin{enumerate}[leftmargin=1.5em]
  \item \textit{Known analogue pairs}: Human-curated pairs of analogous
        theorems as a held-out benchmark.
  \item \textit{Proof reuse prediction}: Whether retrieved theorems help
        prove held-out results.
  \item \textit{Namespace transfer}: Whether the system retrieves useful
        results from different mathematical domains.
  \item \textit{Ablation against baselines}: Compare BM25 over theorem
        names and docstrings; embedding-based search \citep{gao2024leansearch};
        symbol-overlap retrieval \citep{meng2009mepo}; dependency-graph-only
        retrieval; proof-motif-only retrieval; and full multi-layer graph
        retrieval.
  \item \textit{Expert judgment}: Ask Lean users whether retrieved analogies
        are mathematically meaningful.
\end{enumerate}

The key claim to test is: \textit{formal graph structure contains analogy
information not captured by theorem names or text embeddings alone.}

\subsection{Missing Abstraction Detection}

The most ambitious component is a system that detects where mathlib lacks
the right representation.

This turns abstraction discovery into a mathematical refactoring problem
\citep{roycordy2007}. Just as code clone detection identifies repeated code
that should be factored into a shared abstraction, proof clone detection
identifies repeated proof structure that signals a missing lemma, typeclass,
or interface. The analogy with program synthesis library learning is also
direct: both DreamCoder \citep{ellis2021dreamcoder} and Stitch
\citep{bowers2023stitch} identify reusable program substructures by finding
patterns that improve compression; this thesis applies an analogous principle
to formal proofs, exploiting the richer type-theoretic structure of
Lean/mathlib.

\paragraph{Evaluation.}
Evaluation uses three complementary approaches.

\textit{Retrospective}: Use mathlib history. Identify moments when important
abstractions were introduced. Ask whether the system, using only the
pre-abstraction state, would have detected the relevant cluster of
representational pain.

\textit{Prospective}: Run the system on current mathlib, produce suggestions,
and have maintainers judge whether the suggestions are meaningful and whether
accepted suggestions shorten proofs.

\textit{Synthetic}: Remove or hide known lemmas and abstractions from a
subset of the library and test whether the system rediscovers the need for
them.

% -----------------------------------------------------------------------
\section{Technical Architecture}

A plausible architecture has six layers.

\subsection{Layer 1: Trace Collection}
\label{sec:arch-layer1}

Instrument Lean to collect theorem statements, proof terms, tactic traces,
proof states before and after tactics, generated subgoals, used lemmas,
typeclass resolution traces, simplification traces, coercion insertions, and
elaboration metadata. Existing infrastructure from LeanDojo
\citep{yang2023leandojo} and PACT \citep{han2021pact} provides a foundation;
the key extension is the proof-state transition graph and normalization
pipeline.

\subsection{Layer 2: Formal Graph IR}

Convert raw Lean data into a normalized graph intermediate representation
stable across superficial syntax changes. Work by \citet{paliwal2020graph}
and \citet{bansal2019holist} on graph representations for HOL Light provides
relevant prior art; the Lean~4 setting introduces additional structure through
its typeclass and coercion system.

\subsection{Layer 3: Motif and Pattern Mining}

Mine repeated structures including frequent proof-state transitions, theorem
statement schemas, recurring dependency neighborhoods, repeated
coercion/rewrite paths, common proof skeletons, and bottleneck states. This
layer combines symbolic graph algorithms with learned embeddings. The static
dependency analysis of \citet{huch2022} on the Isabelle AFP demonstrates that
large formal library graphs have measurable large-scale structure; this layer
extends that analysis to proof-state dynamics.

\subsection{Layer 4: Retrieval and Analogy}

Build retrieval models over the graph IR using contrastive learning over
multi-layer theorem graphs, with ablations against text-embedding baselines
\citep{gao2024leansearch} and symbol-overlap baselines \citep{meng2009mepo}.
Additional methods to explore include graph kernels, subgraph matching,
nearest-neighbor search over graph embeddings, and hybrid symbolic-neural
retrieval.

\subsection{Layer 5: Abstraction Scoring}

Define metrics for representational inefficiency and abstraction value,
including: repeated motif density, proof compression ratio
\citep{bowers2023stitch}, dependency centrality, namespace-spanning
recurrence, typeclass boundary friction, search entropy, proof-state
branching factor, normalized proof length, and theorem transportability.

\subsection{Layer 6: Interactive Assistant}

Expose the system to users as an assistant for formal mathematics. Given a
proof state, it can suggest analogous proof patterns, useful intermediate
lemmas, repeated hypotheses, alternative transports, and potential missing
abstractions. The final system is not merely an auto-prover; it is a
representation-aware mathematical assistant.

% -----------------------------------------------------------------------
\section{Initial Projects}

\subsection{Project 1: Proof-State Motif Mining in Mathlib}

\textbf{Research question}: Do formal proofs contain recurring structural
motifs that are stable across domains and predictive of proof strategy?

\textbf{Deliverables}: proof-state graph extractor; motif mining algorithm;
motif taxonomy; visualization tool; evaluation on tactic prediction and
theorem retrieval.

\textbf{Paper claim}: Proof-state transition graphs contain reusable
cross-domain motifs not captured by tactic sequences alone.

\subsection{Project 2: Structural Theorem Analogy Retrieval}

\textbf{Research question}: Can formal graph structure retrieve useful
mathematical analogies beyond lexical or embedding-based theorem search?

\textbf{Deliverables}: multi-layer theorem graph representation; retrieval
benchmark with human-curated analogue pairs; comparison against text-embedding
baselines \citep{gao2024leansearch} and symbol-overlap baselines
\citep{meng2009mepo}; case studies of cross-domain analogies.

\textbf{Paper claim}: Multi-layer formal graph retrieval finds structurally
analogous theorems across mathlib and improves proof reuse.

\subsection{Project 3: Representation Inefficiency and Missing Lemma
Discovery}

\textbf{Research question}: Can missing abstractions be detected through
repeated proof patterns and graph-theoretic bottlenecks?

\textbf{Deliverables}: representation inefficiency score; retrospective study
using mathlib history; system-generated missing lemma candidates; evaluation
through proof compression.

\textbf{Paper claim}: Formal libraries contain measurable representational
pain points that can be used to suggest abstractions reducing proof
complexity.

% -----------------------------------------------------------------------
\section{Why This Is Different From Existing AI Theorem Proving}

Most current systems attempt to automate proof search, asking:
\begin{quote}
Given a proof state, what tactic should be applied next?
\end{quote}

This research asks a different question:
\begin{quote}
Given a mathematical landscape, what representation would make proof search
easier?
\end{quote}

Tactic prediction is local; representation discovery is global. Tactic
prediction imitates proof scripts; representation discovery studies why some
proof scripts become short, reusable, and natural. Tactic prediction operates
within the current library; representation discovery asks how the library
itself should evolve. Tactic prediction tries to close goals; representation
discovery tries to invent the objects that make goals close naturally.

This thesis also differs from the closest existing work on formal library
analysis. \citet{li2026networkstructure} extract Mathlib's multilayer
dependency graph and study its macroscopic structure but examine only static
dependencies and do not pursue proof-state dynamics, motif retrieval, or
abstraction detection. \citet{huch2022} applies the same approach to the
Isabelle AFP with the same limitation. \citet{blanchette2015mining} mine the
AFP for descriptive proof statistics but do not use the results operationally.
TacMiner \citep{xin2025tacminer} builds Tactic Dependence Graphs over Lean~4
proofs to discover reusable tactic libraries---the closest prior work to
Project~1---but targets tactic compression rather than representation
discovery. All four works use formal library structure descriptively or for
tactic automation; this thesis uses it as an active substrate for
representation discovery.

The long-term goal is not just an AI that proves existing theorems, but an
AI that helps mathematicians build better mathematical languages.

% -----------------------------------------------------------------------
\section{Why Lean Is the Right Substrate}

Lean is especially suitable for this research because it combines: a large
and actively developed mathematical library \citep{mathlib2020}; expressive
dependent type theory; rich typeclass mechanisms; tactic-based proof
development; machine-checkable proof terms; substantial use of abstraction
and hierarchy; enough library history to study abstraction evolution; an
active community of expert formalizers; and existing extraction infrastructure
\citep{yang2023leandojo, han2021pact}.

The richness of the Lean typeclass system provides additional motivation.
\citet{baanen2022instances} analyzes how mathlib uses Lean's instance
parameter mechanism to organize its algebraic, order, topology, and analysis
hierarchies, showing that representation choices in the typeclass system are
non-trivial, consequential, and frequently interwoven in ways that create
design tension---supporting the thesis's premise that Lean/mathlib contains
substantial traces of representation decisions worth studying.

The HOL Light ecosystem provides a useful comparison point. HOList
\citep{bansal2019holist} and the GNN work of \citet{paliwal2020graph}
demonstrate that graph-structured formal proof data supports meaningful
machine learning. Lean/mathlib offers a richer typeclass hierarchy, a larger
and more systematically organized library, and a more active community.

Other proof assistants are also relevant. The Isabelle AFP has been analyzed
graphically \citep{huch2022} and statistically \citep{blanchette2015mining},
and the Mizar Mathematical Library has a long history as a formal corpus.
These libraries may serve as secondary validation datasets to test whether
mined motifs and detected abstractions generalize across systems.

% -----------------------------------------------------------------------
\section{Expected Contributions}

\paragraph{To AI theorem proving.}
Representation-aware retrieval, proof planning, intermediate lemma discovery,
and abstraction-guided search.

\paragraph{To formal methods.}
Tools for library refactoring, proof maintenance, theorem organization, and
abstraction quality measurement.

\paragraph{To programming languages.}
A treatment of formal mathematics libraries as evolving typed knowledge
systems, with abstraction studied as an empirical software phenomenon.

\paragraph{To machine learning.}
New graph-structured learning problems grounded in rigorous symbolic data.

\paragraph{To mathematics.}
Systems that may help mathematicians discover analogies, transport ideas
across fields, and identify the right objects to define.

\paragraph{To the philosophy and cognitive science of mathematics.}
A computational way to study abstraction, analogy, representation, and proof
compression---operationalizing theoretical frameworks such as structure-mapping
theory \citep{gentner1983} on large-scale formal data \citep{mitchell2021,
zhang2023cognitive}.

% -----------------------------------------------------------------------
\section{Long-Term Vision}

The long-term vision is an AI-native environment for formal mathematics in
which the system does not merely autocomplete proofs, but actively helps
users navigate the space of mathematical representations.

The program synthesis community has shown that library learning from a corpus
of programs is feasible: DreamCoder \citep{ellis2021dreamcoder} and Stitch
\citep{bowers2023stitch} demonstrate that useful abstractions can be
discovered by mining reusable structure and measuring compression. This thesis
pursues an analogous program for formal mathematics, where the corpus is
Lean/mathlib, the abstractions are mathematical definitions and typeclasses,
and correctness is guaranteed by the type checker rather than by test suites.
The richer type-theoretic structure of formal mathematics creates both
additional challenges (the abstraction space includes coercions, universal
properties, and bundled structures with no direct program synthesis analogue)
and additional opportunities (graph structure is richer, and correctness is
mechanically verified).

In this environment, Lean becomes more than a proof checker. It becomes a
microscope for mathematical structure.

The most ambitious version of the thesis is:

\begin{quote}
Formal mathematics libraries are the first large-scale datasets from which
machines can learn the dynamics of mathematical abstraction.
\end{quote}

If this is true, then the future of AI for mathematics is not merely
automated proof. It is automated representation discovery.

The path from the three initial projects to this vision requires
demonstrating, sequentially: (i)~that proof-state motifs exist and are
cross-domain; (ii)~that graph-based analogy retrieval outperforms text-based
retrieval on a held-out benchmark; and (iii)~that the system's
missing-abstraction suggestions correspond to genuine improvements in
mathlib. Each step is independently publishable and provides the empirical
foundation for the next.

% -----------------------------------------------------------------------
\section{Risks and Failure Modes}

\paragraph{Risk 1: Lean artifacts may overwhelm mathematical signal.}
Proof states contain many artifacts from elaboration, typeclass resolution,
coercions, and naming choices. Graph mining may discover Lean-specific noise
rather than mathematical structure.
\textit{Mitigation}: normalization, quotienting, ablation studies,
cross-namespace validation, and comparison across different formalizations.

\paragraph{Risk 2: Graph similarity may not correspond to meaningful analogy.}
Two theorem graphs may look similar for superficial reasons, while deep
analogies may be structurally distant.
\textit{Mitigation}: combine symbolic graph features with expert judgment,
proof reuse evaluation, and retrieval-based downstream tasks. Benchmark
against human-curated analogue pairs.

\paragraph{Risk 3: Missing abstraction detection may produce vague suggestions.}
It is easy to identify repeated patterns; it is harder to propose a clean
abstraction.
\textit{Mitigation}: begin with missing lemma discovery and proof compression
before attempting new structure invention.

\paragraph{Risk 4: Full proof-state extraction may be technically expensive.}
Tracing all of mathlib may be large, brittle, or slow. LeanDojo
\citep{yang2023leandojo} provides starting infrastructure but is not
optimized for full proof-state transition graph extraction.
\textit{Mitigation}: start with selected domains, cache aggressively, support
incremental extraction, and focus first on tactic-level traces and dependency
graphs.

\paragraph{Risk 5: The project may become too philosophical.}
The big picture is attractive, but the field will only move if concrete
systems and evaluations exist.
\textit{Mitigation}: anchor the research agenda in three initial artifacts
(motif miner, analogy retriever, missing abstraction detector), each with a
concrete, reproducible evaluation.

% -----------------------------------------------------------------------
\section{Positioning Statement}

This research direction is not only ``AI for theorem proving,'' ``graph
mining on mathlib,'' or ``retrieval for Lean.'' It is the study and
automation of how formal mathematical representations are created,
compressed, reused, and transported.

\begin{quote}
We use Lean/mathlib as a machine-readable record of mathematical abstraction.
By extracting and analyzing proof-state graphs, theorem dependency graphs,
and typeclass structures, we build systems that discover proof motifs,
retrieve cross-domain analogies, detect missing abstractions, and guide
representation-aware proof search.
\end{quote}

Shorter: \textit{this project turns formal mathematics libraries into datasets
for learning mathematical representation discovery.}

% -----------------------------------------------------------------------
\section{First Paper Sketch}

\paragraph{Possible title.}
\textit{Proof Motifs in Formal Mathematics: Mining Reusable Proof-State
Patterns from Lean/mathlib}

\paragraph{Core claim.}
Formal proof traces contain recurring structural motifs that appear across
mathematical domains and improve theorem retrieval and proof-step prediction.

\paragraph{Contributions.}
\begin{enumerate}[leftmargin=1.5em]
  \item A proof-state graph extraction pipeline for Lean~4, extending
        LeanDojo \citep{yang2023leandojo} with transition-graph structure.
  \item A normalized representation of tactic-induced proof-state transitions.
  \item A motif mining algorithm for formal proofs.
  \item An empirical taxonomy of recurring proof motifs in mathlib.
  \item Evidence that motif features improve theorem retrieval or proof
        guidance.
\end{enumerate}

\paragraph{Evaluation.}
Motif recurrence across namespaces; prediction of next tactic category;
retrieval of proof-similar theorems; human evaluation of motif
interpretability; ablation against tactic-sequence baselines, theorem-text
baselines \citep{gao2024leansearch}, and dependency-graph-only baselines
\citep{huch2022}.

\paragraph{Why this paper matters.}
It demonstrates that formal proof traces contain reusable higher-level
structure beyond tactic sequences. This justifies the broader agenda of
representation-aware AI for mathematics and distinguishes this work from
tactic prediction systems \citep{lample2022hypertree, yang2023leandojo} and
prior static library analysis \citep{huch2022, blanchette2015mining}.

% -----------------------------------------------------------------------
\section{Final Thesis}

The history of mathematics is partly the history of better representations.
New objects, abstractions, and languages make previously impossible proofs
natural. Until recently, this process was visible only through books, papers,
lectures, and human interpretation. Formal mathematics changes that. It
records mathematical representation in executable, machine-readable form.

Lean/mathlib \citep{mathlib2020} therefore gives us a new scientific object:
a large-scale formal trace of mathematical abstraction in action.

The opportunity is to build systems that can read this trace---not merely to
imitate proof steps, but to learn the structural patterns by which mathematics
organizes itself.

\begin{quote}
By treating formal mathematics libraries as graph-structured records of
representation, proof, and abstraction, we can build AI systems that retrieve
deep analogies, discover reusable proof motifs, identify missing abstractions,
and guide mathematicians toward better representations. This shifts AI for
mathematics from local proof search toward representation discovery, the level
at which much of mathematical creativity actually occurs.
\end{quote}

This is the direction worth leading.

% -----------------------------------------------------------------------
\bibliographystyle{plainnat}
\bibliography{references}

\end{document}
